\section{Introduction}\label{sec:intro}
Yoga is practised worldwide and is increasingly learned alone, from videos and phone apps; a systematic review of published case reports documented adverse events associated with yoga \cite{cramer2013yoga}. A camera-based coach that tells a practitioner which joint is off, in their own language, is therefore attractive, and real-time body-landmark estimators that run on a phone \cite{bazarevsky2020blazepose,lugaresi2019mediapipe} make it technically possible. A wave of 2025--2026 systems has followed: sensor-augmented and IoT designs \cite{yuan2025gtanet,sathyapriya2026posture,kamble2026yogamat,saravanapandian2025mat}, a virtual-reality coach \cite{srinivasakumar2025yogavr}, edge and hybrid pipelines \cite{gadhvi2025posepilot,haldar2026framework,vallabhaneni2026hybrid}, joint-level correction from a single image \cite{chi2026sage}, an accessibility-oriented system \cite{aggarwal2026chair}, and a clinical validation of on-device pose estimation for resistance training \cite{heo2026ondevice}.

\runin{Open problems.} Reading this literature, we identify six problems that remain open (Table~\ref{tab:related} summarises what each system reports).
\begin{itemize}\setlength{\itemsep}{1pt}\setlength{\topsep}{2pt}
\item (G1) Evaluation. Reviews note heterogeneous metrics and uncertain generalisability \cite{ozsezer2025review,koroglu2026review}, and typical results are measured on a random split of a small image set (for example 96.09\% on a 256-image test split of 1{,}280 images \cite{mohiuddin2025skeleton}); random splits place near-identical frames and photos on both sides and are a known source of inflated results \cite{kapoor2023leakage}. None of the abstracts we read reports how often a pose outside the vocabulary is wrongly named.
\item (G2) Temporal behaviour. Pose transition is described as under-explored \cite{dittakavi2025pose2pose}; recurrent models appear in several systems \cite{gadhvi2025posepilot,haldar2026framework}, but their abstracts report no evaluation on videos of unseen people.
\item (G3) Personalisation. Where recent yoga systems personalise, they rely on extra sensors \cite{vallabhaneni2026hybrid,saravanapandian2025mat}.
\item (G4) Hidden joints. Occlusion is a stated difficulty of the domain \cite{verma2020yoga82}, yet the yoga systems above do not describe in their abstracts how a hidden joint is scored.
\item (G5) Safety of generated advice. Vision--language coaching still shows substantial gaps against human coaches \cite{zuo2025formcoach}, and a system that uses a language-model voice assistant \cite{srinivasakumar2025yogavr} reports no screening of its output in its abstract.
\item (G6) Deployment. Systems need mats, headsets, robots or wearables \cite{kamble2026yogamat,saravanapandian2025mat,srinivasakumar2025yogavr,sathyapriya2026posture}, and none of the abstracts we read reports guidance beyond English.
\end{itemize}

\runin{Approach and research questions.} We propose AsanaAI, a coach designed around these problems and evaluated only on data its models never trained on. We ask:
\begin{itemize}\setlength{\itemsep}{1pt}\setlength{\topsep}{2pt}
\item RQ1: Does a name-and-veto cascade of small residual MLPs name poses accurately on held-out photos while rarely naming a pose that is not in its vocabulary, compared with standard classifiers?
\item RQ2: Does a skeleton-graph temporal network recognise held poses and separate them from transitions on videos of unseen people, what input rate does it need, and does the skeleton graph beat simpler sequence models?
\item RQ3: Is declining to score hidden joints preferable to recovering them by mirroring or by mesh recovery?
\item RQ4: Does a personal range-of-motion profile change feedback only where intended, and can the system run on commodity hardware with acceptable latency?
\end{itemize}

\runin{Contributions.}
\begin{enumerate}
\item \emph{Cue-verified supervision:} instructor videos are turned into training labels only where a speech transcript, a biomechanical rule engine and a steadiness test all agree.
\item \emph{A name-and-veto cascade} of two 3-head residual MLPs with an explicit ``not one of mine'' answer, and a controlled comparison with five standard classifiers on two held-out photo sets, including a result that goes against our own design (Section~\ref{sec:rq1}).
\item \emph{A skeleton-graph flow check} that separates held poses from transitions, with its input-rate requirement measured and a controlled comparison with four simpler sequence models and a graph-free variant, which shows parity, not superiority.
\item \emph{Occlusion-aware scoring} that abstains on hidden joints, compared with mirroring and mesh recovery.
\item \emph{Personal calibration and safety-screened guidance} in English, Hindi and Bengali, with the image never leaving the device.
\item \emph{Evidence on commodity hardware:} server and phone measurements and an end-to-end test on a mid-range phone, with live screenshots of the deployed interface.
\end{enumerate}
